\chapter[The Pattern Describing Itself]{The Pattern\\Describing Itself}

\vspace{1em}
\hfill
\begin{minipage}[c]{0.6\textwidth}
\raggedleft
\small\itshape
The growing good of the world is partly dependent on unhistoric acts; and that things are not so ill with you and me as they might have been, is half owing to the number who lived faithfully a hidden life, and rest in unvisited tombs.\\[0.5em]
\upshape\footnotesize George Eliot, \textit{Middlemarch} (1872)
\end{minipage}%
\hspace{0.5em}
\begin{minipage}[c]{0.12\textwidth}
% \includegraphics[width=\textwidth]{figures/ch15_closing.png}
\end{minipage}
\vspace{1.5em}

\noindent \textsc{This book is a finite pattern of symbols.} A program can reproduce it, so it is one of the Library's describable outputs. Under Picture A, that fact can be read ontologically; under Picture B, it is enough that this copy exists here and that you are reading it now.

So the book turns, at the end, and looks at itself. It has spent fourteen chapters saying the structure is indifferent. Here is the last thing it can say, and it is not an argument.

Gather the cold first. $M$ weights candidate programs, relative to a machine, but supplies no observer-occurrence measure and selects no system as real. The postulate denies a view-independent decomposition while leaving Picture A and Picture B open. Everett is optional. No theorem here explains why these laws obtain, why there is anything at all, or how a first-person index should be assigned. A greater mind, if one exists, would still have a view rather than the view from nowhere. The room is as cold as the book can honestly make it, without turning an open question into a fact.

And here you are.

In a local situation, finite, with this book in your hands and, very likely, people in the next room or the next call or the unanswered message. No global measure turns those relations into obligations for you. Their warmth is made between finite people through attention and action. It is the cut we keep making toward each other, the one the formalism will not make for us.

This is the move the book was for. Care is not a conclusion the cosmology provides. It is something a person does here, with limited causal reach. The formalism does not require or reward it. A situation with care is nevertheless not the same for the people in it as a situation without care.

Consider the object itself: paper, or a lit screen, a finite arrangement of symbols. Someone set these words down with some care, to reach a reader they might never meet. You are reading them. That is all this is, and it is not small. A person described a pattern and chose, in the describing, to be kind about it. Reading carefully can answer from the other side.

Rage is one honest response to the cold: rage at unanswered questions, brute facts, and the absence of guaranteed rescue. It may also pass. What can remain is neither despair nor comfort, but a finite person deciding what kind of presence to be with the time available.

Camus, at the end of his own cold argument, wrote \emph{il faut imaginer Sisyphe heureux}: one must imagine Sisyphus happy. The \emph{must} is not a derivation. It is a commitment made without waiting for a cosmology to authorize it. The author wrote this book as that commitment, performed. You may read it with care as another performance. Whether anything greater sees it is not a question this book has settled. We see it.

Is it enough? I do not know. I cannot reach every place where tending failed, and I cannot promise that the braid you can hold answers suffering beyond your reach. I hope it does. The hope is not a proof, and I will not hand it to you as one. But choosing does not wait on the answer. We do it unsure, which is the only way it was ever going to be done.

\vspace{1.2em}
\noindent Stay.
